\section*{Part X: Bayes Institute: From Rented Space to a Credible Campus}
\addcontentsline{toc}{section}{Part X: Bayes Institute: From Rented Space to a Credible Campus}


\subsection{70. The starting state}

Bayes has not established a physical campus or a proven student volume justifying a permanent teaching center. The founder has technical teaching and AI expertise but limited weekly time alongside employment. Land, offices and GPUs should therefore be considered future \textbf{options tied to evidence}, not founding prerequisites.

Do not mix prospective personal residential construction with Bayes company property or represent a planned studio as a built institute. Use formally documented, lawful business arrangements for any shared space, with tax, safety, insurance and related-party implications reviewed.


\subsection{71. Phase 0: temporary access}

Use rented classrooms for demonstrations, partner laboratories for specific assignments, and occasional in-person capstones. Measure attendees, perceived benefits, independent learning, setup time, facilities costs, and recurring demand. A suitable venue is an experiment in educational value and community, not a prototype campus merely because learners like meeting.

The goal is to discover which physical moments create the most value and whether the value is worth recurring fixed cost.


\subsection{72. Phase 1: the small studio decision}

A modest leased or appropriately authorized shared studio may become rational if it materially improves teaching evidence, recording, staff collaboration or recurring learner attendance. Compare monthly full costs with continued pay-per-use space. Establish a named operating owner, booking system, safety rules, learner accessibility and realistic utilization. Do not lease more area because the institute may be ``big soon.''

A 3- or 6-month pilot with an exit right can be more valuable than a 5-year lease at a superficially low rate.


\subsection{73. Phase 2: a specialized applied lab}

A partner-college AI/robotics/engineering lab could provide equipment and local access while Bayes supplies curriculum and research supervision. Ensure intellectual property, funding, student safeguarding, equipment ownership, warranties, access schedule, branding and publication rights are specified. If the partner funds capex, the institute still has a continuing management and quality obligation.

Build an owned lab only after actual workloads and research questions justify facility specifications and replacement reserves.


\subsection{74. Phase 3: site-search charter for a permanent academic nucleus}

Define a 10-20-year academic question: which students and faculty, what STEM disciplines, what research intensity, what residential requirement, and what legal path? Commission alternative region analyses, then detailed micro-market visits and parcel due diligence. Establish an independent adviser panel spanning academic leadership, property law, planning, engineering and finance.

The output is a \textbf{location decision memorandum with hard gates}, not a broker's short list.


\subsection{75. A Bayes location comparison to research, not a recommendation}

A reasonable initial longlist may include Pune outskirts for technical talent and academic/industry proximity; selected Mumbai-region sites for employer and international access but potentially higher prices; Nashik and Nagpur for alternative land and urban ecosystems; and other regions if funding or research partners create an advantage. Evaluate actual travel times, faculty family needs, water/power, planning rules and cost per lawful built area.

An option to lease an urban research nucleus and defer greenfield land can dominate the urge to identify 1 forever-campus location immediately.


\subsection{76. The university regulatory gate is separate from the land gate}

UGC distinguishes lawful degree-awarding universities from ordinary education providers and emphasizes that the relevant state controls criteria for state private universities [P2]. Maharashtra's 2023 private-university statute and 2024 rules form a relevant framework, subject to amendments and the chosen institutional structure [P3]. An official 2024 rules document describes a project application with a detailed project report and a specified processing fee; verify the current fee and all sponsor, land, endowment, approval and campus conditions before any material expense [P3].

A plot that is perfect for a non-degree training studio is not automatically a valid site under a particular future university establishment route.


\subsection{77. Phase 4: finance a minimum viable campus, not an empty master plan}

Define a 1st academic phase that can lawfully operate and meet the intended mission: perhaps a few teaching/research buildings, essential services, and appropriately staged student/faculty support. Construct only once funding for phase-1 capex, contingencies, opening deficits and ongoing reserves is credible. Donor pledges must be converted into cash timing; project loans need realistic repayment capacity; restricted assets must be ring-fenced.

The 1st phase should be useful even if future phases take 10 additional years or never occur.


\subsection{78. A Bayes physical-expansion decision tree}

If the educational value is intermittent, \textbf{book access}. If regular use is proven but future geography uncertain, \textbf{lease with exit protection}. If specialized research needs are concentrated and partners are credible, \textbf{co-invest or share labs}. If the institute has recurring physical requirements, stable operators and cash, \textbf{consider a small owned nucleus}. If degree-scale land requirements, academic governance, sponsor and multi-year financing are established, \textbf{evaluate a permanent campus}. At each transition, compare the option to doing nothing and strengthening the core.

This is a sequencing framework, not a prediction that Bayes will eventually own land.
